\chapter{Conclusion}
\label{sec:conclusion}
% Answer each research question and address how the limitations of the study qualify the conclusion.
%Write your conclusion here. Be sure that the relation between the research gap and your contribution is clear. Be honest about how limitations in the study qualify the answer on the research question.


This thesis addressed a gap in the tool-use literature: the assumption that
tools are given. We asked whether language models can synthesize functional
API tool servers from documentation, whether those servers enable effective
agent task completion, and what role parametric knowledge plays in both.
Below we answer each research question and qualify the answers with the
relevant limitations of the study.

\paragraph{RQ1: Can LLMs synthesize functional API tool servers from documentation alone?}
Yes, for most APIs tested. Eleven of fourteen APIs achieve SYNTH $\geq 75\%$ in the docs condition, and the frontier-tier median PASS of 83\% demonstrates these servers are not just syntactically valid but produce real task-completion utility in practice. The dominant failure mode is coverage gap (omitted endpoints) rather than structural failure (broken or non-executable servers). Notably, successful synthesis does not require frontier-scale models for all APIs: mid-tier models achieve competitive PASS on simpler, high-familiarity APIs (e.g., GitHub, Mastodon), though frontier models produce more reliable and consistent coverage across the full benchmark. This answer is qualified by generalizability: all APIs in the benchmark are REST-based with English documentation, and synthesis feasibility for larger or differently structured APIs remains to be established.

\paragraph{RQ2: Do synthesized servers enable effective agent task completion, compared to human-authored references?}
Yes, at a frontier-tier median PASS of 83\% in the docs condition, and synthesized servers match or exceed the two available human-authored references. On GitHub, no-docs synthesis achieves 83\% PASS (81--90\% across variance-study runs) against an official server baseline of 65\%; on Zulip, docs-condition synthesis scores 80\% against a baseline of 68\%. The qualification in this case is that agent performance is partially
confounded by the fixed evaluation agent (Gemini~2.5~Pro); agent-limited failures identified in our benchmark may not reflect the same limitations with a different evaluator.

\paragraph{RQ3: To what extent do synthesizers rely on parametric knowledge, and does this affect performance?}
Substantially, and in a way that is API-dependent. High-familiarity APIs show near-zero HTTP mutation adherence, indicating synthesis models write almost entirely from memory. The performance consequence is non-monotonic: parametric reliance helps on GitHub (where model memory is accurate and complete) but does not on Stripe (where documentation provides additional coverage). This answer is limited by the fact that adherence is measured for a fixed set of 4--6 mutated parameters per API; the pattern may differ for other models or mutation strategies.

\paragraph{RQ4: How does documentation availability affect synthesis and agent performance?}
The aggregate effect is smaller than expected: the no-docs condition median (88\%) meets or exceeds the docs condition median (83\%), driven by high-familiarity APIs where model memory substitutes for documentation. Documentation helps for less familiar or more complex APIs (Stripe, Zulip) and hurts for highly familiar or simply-structured ones (GitHub, Mastodon), suggesting the effect is mediated by training coverage rather than documentation quality per se. This finding
is subject to the static-snapshot limitation: documentation was scraped at a single point in time, and results may differ for documentation that is more consistently current or higher quality.

\paragraph{RQ5: Does synthesis quality predict agent task performance?}
Strongly: $r = 0.937$ ($p < 0.0001$, $n = 209$ across all models). SYNTH, assessable without running the evaluation agent, is a reliable proxy for PASS. This holds across APIs and conditions, with the main exceptions being cases where the
evaluation agent compensates for synthesis gaps using its own domain knowledge (Spoonacular) or where the agent fails despite sufficient tools (agent-limited cases). The correlation is measured on our 14-API benchmark; whether it holds at this strength for a wider or different API distribution is unknown.

\section*{Research Gap and Contribution}

Prior work treats the tool layer as given. We have shown that synthesizing it automatically is feasible today for most practically relevant REST APIs, and that the primary bottleneck for large APIs is context-window coverage of documentation rather than a fundamental capability limit. The key novelty is the synthesis \emph{target}: a complete, protocol-compliant MCP server rather than individual tool stubs, which is what enables genuine end-to-end evaluation through live API execution. The SYNTH metric and the controlled ablation design --- isolating parametric knowledge from documentation grounding --- are methodological contributions that generalise beyond this benchmark to any evaluation of documentation-grounded
code generation. A preliminary mechanistic analysis (Chapter~\ref{sec:mechanistic}) grounds the parametric gravity phenomenon in transformer representations, connecting behavioural findings to the interpretability literature.

\section*{Future Work}

Several directions emerge directly from the findings.

\textbf{Coverage scaling.} The most immediate target is decomposed synthesis for large APIs (Jira, eBay Sell): synthesizing one resource group at a time and assembling the result addresses the context-window bottleneck without requiring larger models.

\textbf{Proprietary and non-public APIs.} All 14 benchmark APIs have public documentation. Extending synthesis to internal or proprietary APIs --- where documentation may be incomplete, inconsistently maintained, or confidential --- introduces new challenges around grounding signal quality and access control that merit separate investigation.

\textbf{Neuro-symbolic and constraint-driven synthesis.} A natural follow-on to the full-server setting is to enforce protocol constraints symbolically at synthesis time: type-checking tool schemas, validating endpoint paths against an API registry, or running a lightweight stub call before committing to an implementation. Such neuro-symbolic approaches could eliminate the most common \verdict{tool\_schema} failures without requiring re-synthesis.

\textbf{Adaptive documentation inclusion.} Our findings suggest that the optimal synthesis context is API-specific and predictable from familiarity. A practical system could automate this decision: a pre-synthesis probe (or mutation adherence mini-run) could estimate model confidence in the API's parametric representation and gate documentation inclusion accordingly.

\textbf{Specialised multi-agent synthesis pipelines.} The single-agent synthesis scaffold used here is minimal by design. More capable pipelines could include a planner that decomposes synthesis by resource group, a verifier that tests generated tools against the live API, and an assembler that merges partial servers --- analogous to specialised coding agents for software engineering tasks.

\textbf{Live versioning benchmarks.} Extending the probing analysis to frontier-scale open-weight models would test whether the layer-transition pattern generalises. A live benchmark that re-runs synthesis as APIs evolve would capture the dynamic challenge that static snapshots cannot, particularly for APIs at the Stripe/Zulip familiarity tier where parametric memory and documentation are in active tension.
